\documentclass[../../main.tex]{subfiles}

\begin{document}

\section{完备度量空间}\label{section:ch09-完备度量空间}

度量空间自身的结构太贫瘠而不能在数学分析的有趣问题的研究中富有成果. 然而, 值得注意的是, 通过考虑仅具有一条额外性质即完备性的度量空间, 我们便能够建立大量有趣而重要的结果. 我们将下一章贡献给完备度量空间的三个基本定理.

\begin{definition}[Cauchy 序列与完备性]\label{definition:ch09-Cauchy完备}
度量空间 $(X,\rho)$ 中的序列 $\{x_{n}\}$ 称为是 Cauchy 序列, 若对每个 $\varepsilon>0$, 存在指标 $N$ 使得
\begin{align*}
\text{若 }n,m\geqslant N,\quad\text{则 }\rho(x_{n},x_{m})<\varepsilon.
\end{align*}

度量空间 $X$ 称为完备的, 若 $X$ 中的每个 Cauchy 序列收敛到 $X$ 中的点.
\end{definition}

这推广了我们首先在第 1 章对于实数序列, 接着在第 7 章对于赋范线性空间序列考虑过的 Cauchy 序列的概念. 对于一般的度量空间, 如同赋范线性空间的情形, 一个收敛的序列是 Cauchy 的, 而一个 Cauchy 序列是有界的 (见习题 38).

实数的完备性公理等价于度量空间 $\mathbf{R}$ 的完备性. 由此我们可推断出每个欧氏空间 $\mathbf{R}^{n}$ 是完备的. 此外, 7.3 节证明的 Riesz-Fischer 定理告诉我们, 对于实数的 Lebesgue 可测子集 $E$ 和 $1\leqslant p\leqslant\infty$, $L^{p}(E)$ 是完备的.

\begin{proposition}\label{proposition:ch09-Cab完备}
令 $[a,b]$ 为实数的闭有界区间. 则具有由最大值范数诱导的度量的 $C[a,b]$ 是完备的.
\end{proposition}

\begin{proof}
令 $\{f_{n}\}$ 为 $C[a,b]$ 中的 Cauchy 序列. 首先假定存在收敛的级数 $\sum_{k=1}^{\infty}a_{k}$ 使得
\begin{align}
\text{对所有 }k,\|f_{k+1}-f_{k}\|_{\max}\leqslant a_{k}.\label{align:ch09-Cab快速Cauchy}
\end{align}
由于
\begin{align*}
\text{对所有 }n,k,\quad f_{n+k}-f_{n}=\sum_{j=n}^{n+k-1}[f_{j+1}-f_{j}],
\end{align*}
\begin{align*}
\text{对所有 }n,k,\|f_{n+k}-f_{n}\|_{\max}\leqslant\sum_{j=n}^{n+k-1}\|f_{j+1}-f_{j}\|_{\max}\leqslant\sum_{j=n}^{\infty}a_{j}.
\end{align*}
令 $x$ 属于 $[a,b]$. 则
\begin{align}
\text{对所有 }n,k,\quad|f_{n+k}(x)-f_{n}(x)|\leqslant\sum_{j=n}^{\infty}a_{j}.\label{align:ch09-Cab逐点估计}
\end{align}
级数 $\sum_{k=1}^{\infty}a_{k}$ 收敛, 因此 $\{f_{n}(x)\}$ 是实数的 Cauchy 序列. 实数集是完备的. 记 $\{f_{n}(x)\}$ 的极限为 $f(x)$. 在\eqref{align:ch09-Cab逐点估计}中取 $k\to\infty$ 的极限, 得出
\begin{align*}
\text{对所有 }n\text{ 和所有 }x\in[a,b],\quad|f(x)-f_{n}(x)|\leqslant\sum_{j=n}^{\infty}a_{j}.
\end{align*}
我们从该估计式推出 $\{f_{n}\}$ 在 $[a,b]$ 上一致收敛于 $f$. 由于每个 $f_{n}$ 是连续的, $f$ 也如此. 通过注意到若一个 Cauchy 序列有收敛的子序列则它收敛, 而 $C[a,b]$ 中的每个 Cauchy 序列有一个使得\eqref{align:ch09-Cab快速Cauchy}成立的子序列, 一般的情形可从特殊的情形得到.
\end{proof}

一般来说, 完备度量空间的子空间是不完备的. 例如, 实数的开有界区间是不完备的, 而 $\mathbf{R}$ 是完备的. 然而, 对那些完备的子空间有以下简单的刻画.

\begin{proposition}\label{proposition:ch09-完备子空间闭集}
令 $E$ 为完备度量空间 $X$ 的紧子集. 则度量子空间 $E$ 是完备的当且仅当 $E$ 是 $X$ 的闭子集.
\end{proposition}

\begin{proof}
首先假定 $E$ 是 $X$ 的闭子集. 令 $\{x_{n}\}$ 为 $E$ 中的 Cauchy 序列. 则 $\{x_{n}\}$ 可被考虑为 $X$ 的 Cauchy 序列而 $X$ 是完备的. 因此 $\{x_{n}\}$ 收敛到 $X$ 中的点 $x$. 根据命题 6, 由于 $E$ 是 $X$ 的闭子集, $E$ 中的收敛序列的极限属于 $E$. 于是 $x$ 属于 $E$, 因此 $E$ 是完备度量空间.

为证明逆命题, 假定 $E$ 是完备的. 根据命题 6, 为证明 $E$ 是 $X$ 的闭子集, 我们必须证明 $E$ 中的收敛的序列的极限也属于 $E$. 令 $\{x_{n}\}$ 为 $E$ 中收敛到 $x\in X$ 的序列. 但一个收敛的序列是 Cauchy 的. 因此根据 $E$ 的完备性, $\{x_{n}\}$ 收敛到 $E$ 中的点. 但度量空间的收敛序列仅有一个极限. 因此 $x$ 属于 $E$.
\end{proof}

\begin{theorem}\label{theorem:ch09-完备空间例子}
以下是完备度量空间:
\begin{enumerate}[(i)]
\item 欧氏空间 $\mathbf{R}^{n}$ 的每个非空闭子集.
\item 对实数的可测子集 $E$ 和 $1\leqslant p\leqslant\infty$, $L^{p}(E)$ 的每个非空闭子集.
\item $C[a,b]$ 的每个非空闭子集.
\end{enumerate}
\end{theorem}

\begin{definition}[直径与收缩序列]\label{definition:ch09-直径收缩}
对于度量空间 $(X,\rho)$ 的非空子集 $E$, 我们定义 $E$ 的直径 $\mathrm{diam}E$ 为
\begin{align*}
\mathrm{diam}E=\sup\{\rho(x,y)\mid x,y\in E\}.
\end{align*}
我们说 $E$ 是有界的, 若它有有限的直径. $X$ 的非空子集的递减序列 $\{E_{n}\}_{n=1}^{\infty}$ 称为收缩序列, 若
\begin{align*}
\lim_{n\to\infty}\mathrm{diam}(E_{n})=0.
\end{align*}
\end{definition}

第 1 章的集套定理告诉我们实数的非空闭集的收缩序列的交由单点组成. 这可推广如下.

\begin{theorem}[Cantor 交定理]\label{theorem:ch09-Cantor交}
令 $X$ 为度量空间. 则 $X$ 是完备的当且仅当只要 $\{F_{n}\}_{n=1}^{\infty}$ 是 $X$ 的非空闭子集的收缩序列, 就存在点 $x\in X$ 使得 $\bigcap_{n=1}^{\infty}F_{n}=\{x\}$.
\end{theorem}

\begin{proof}
首先假设 $X$ 是完备的. 令 $\{F_{n}\}_{n=1}^{\infty}$ 为 $X$ 的非空闭子集的收缩序列. 对每个指标 $n$, 选取 $x_{n}\in F_{n}$. 我们宣称 $\{x_{n}\}$ 是一个 Cauchy 序列. 事实上, 令 $\varepsilon>0$. 存在指标 $N$ 使得 $\mathrm{diam}F_{N}<\varepsilon$. 由于 $\{F_{n}\}_{n=1}^{\infty}$ 是下降的, 若 $n$, $m\geqslant N$, 则 $x_{n}$ 和 $x_{m}$ 属于 $F_{N}$, 因此 $\rho(x_{n},x_{m})\leqslant\mathrm{diam}F_{N}<\varepsilon$. 于是 $\{x_{n}\}$ 是一个 Cauchy 序列. 由于 $X$ 是完备的, 该序列收敛到某个 $x\in X$. 然而, 对每个指标 $n$, $F_{n}$ 是闭的且对于 $k\geqslant n$, $x_{k}\in F_{n}$, 于是 $x$ 属于 $F_{n}$. 因此 $x$ 属于 $\bigcap_{n=1}^{\infty}F_{n}$. 这个交集不可能含有两个点, 否则 $\lim_{n\to\infty}\mathrm{diam}F_{n}\neq 0$.

为证明逆命题, 假定对任何 $X$ 的非空闭子集的收缩序列 $\{F_{n}\}_{n=1}^{\infty}$, 存在点 $x\in X$ 使得 $\bigcap_{n=1}^{\infty}F_{n}=\{x\}$. 令 $\{x_{n}\}$ 为 $X$ 中的 Cauchy 序列. 对每个指标 $n$ 定义 $F_{n}$ 为非空集 $\{x_{k}\mid k\geqslant n\}$ 的闭包. 则 $\{F_{n}\}$ 是非空闭集的下降序列. 由于 $\{x_{n}\}$ 是 Cauchy 的, 序列 $\{F_{n}\}$ 是收缩的. 因此, 根据假设, 存在 $X$ 中的点 $x$ 使得 $\{x\}=\bigcap_{n=1}^{\infty}F_{n}$. 对每个指标 $n$, $x$ 是 $\{x_{k}\mid k\geqslant n\}$ 的闭包点, 因此任何中心在 $x$ 的球与 $\{x_{k}\mid k\geqslant n\}$ 有非空交. 因此我们可以归纳地选取一个严格递增的自然数序列 $\{n_{k}\}$, 使得对每个指标 $k$, $\rho(x,x_{n_{k}})<1/k$. 子序列 $\{x_{n_{k}}\}$ 收敛到 $x$. 由于 $\{x_{n}\}$ 是 Cauchy 的, 整个序列 $\{x_{n}\}$ 收敛到 $x$ (见习题 39). 因此 $X$ 是完备的.
\end{proof}

Cantor 交定理的一个非常粗略的几何解释是度量空间不完备, 因为它有"洞". 若 $X$ 是不完备的度量空间, 它总是可以最小限度地适当放大而成为完备的. 例如, 有理数集是不完备的, 但它是完备空间 $\mathbf{R}$ 的稠密度量子空间. 作为进一步的例子, 令 $X=C[a,b]$, 现在考虑范数 $\|\cdot\|_{1}$, 它承袭自 $L^{1}[a,b]$. 度量空间 $(X,\rho_{1})$ 是不完备的. 但它是完备度量空间 $L^{1}[a,b]$ 的稠密度量子空间. 这两个具体例子有十分抽象的推广意义. 我们在习题 50 中概括以下定理的证明.

\begin{theorem}[完备化]\label{theorem:ch09-完备化}
令 $(X,\rho)$ 为度量空间. 则存在一个完备度量空间 $(\widetilde{X},\widetilde{\rho})$, 使得 $X$ 是 $\widetilde{X}$ 的稠密子集且
\begin{align*}
\text{对所有 }u,v\in X,\rho(u,v)=\widetilde{\rho}(u,v).
\end{align*}
\end{theorem}

我们称以上描述的度量空间为 $(X,\rho)$ 的完备化. 在度量空间的背景下完备化是唯一的. 任何两个完备化在 $X$ 上的恒等映射的等距关系下是等距的.

\begin{exercise}\label{exercise:ch09-38}
在度量空间 $X$ 中, 证明: (i) 收敛的序列是 Cauchy 的; (ii) 一个 Cauchy 序列是有界的.
\end{exercise}

\begin{exercise}\label{exercise:ch09-39}
在度量空间 $X$ 中, 证明一个 Cauchy 序列收敛当且仅当它有一个收敛的子序列.
\end{exercise}

\begin{exercise}\label{exercise:ch09-40}
假定 $\{x_{n}\}$ 是完备度量空间 $(X,\rho)$ 的序列且对每个 $n$, $\rho(x_{n},x_{n+1})<1/2^{n}$. 证明 $\{x_{n}\}$ 收敛. 若对每个指标 $n$, $\rho(x_{n},x_{n+1})<1/n$, $\{x_{n}\}$ 收敛吗?
\end{exercise}

\begin{exercise}\label{exercise:ch09-41}
提供一个实数的非空闭集的下降可数族的例子, 其中实数的交集是空集. 这与 Cantor 交定理矛盾吗?
\end{exercise}

\begin{exercise}\label{exercise:ch09-42}
令 $\rho$ 和 $\sigma$ 为非空集合 $X$ 上的等价度量. 证明 $(X,\rho)$ 是完备的当且仅当 $(X,\sigma)$ 是完备的.
\end{exercise}

\begin{exercise}\label{exercise:ch09-43}
证明两个完备度量空间的乘积是完备的.
\end{exercise}

\begin{exercise}\label{exercise:ch09-44}
对于从度量空间 $(X,\rho)$ 到度量空间 $(Y,\sigma)$ 的映射 $f$, 证明 $f$ 是一致连续的当且仅当对 $X$ 中的任何两个序列 $\{u_{n}\}$ 和 $\{v_{n}\}$,
\begin{align*}
\text{若 }\lim_{n\to\infty}\rho(u_{n},v_{n})=0,\quad\text{则 }\lim_{n\to\infty}\sigma(f(u_{n}),f(v_{n}))=0.
\end{align*}
\end{exercise}

\begin{exercise}\label{exercise:ch09-45}
用以下概要证明一致连续映射的延拓性质: 令 $X$ 和 $Y$ 为度量空间, 其中 $Y$ 是完备的, 而 $f$ 是从 $X$ 的子集 $E$ 到 $Y$ 的一致连续映射. 则 $f$ 有唯一的一致连续延拓 $\overline{f}$, $\overline{f}$ 是从 $\overline{E}$ 到 $Y$ 的一致连续映射.
\begin{enumerate}[(i)]
\item 证明 $f$ 将 $E$ 中的 Cauchy 序列映到 $Y$ 中的 Cauchy 序列.
\item 对 $x\in\overline{E}$, 选取 $E$ 中收敛到 $x$ 的序列 $\{x_{n}\}$ 且定义 $\overline{f}(x)$ 为 $\{f(x_{n})\}$ 的极限. 用习题 44 证明 $\overline{f}(x)$ 是恰当定义的.
\item 证明 $\overline{f}$ 在 $\overline{E}$ 上是一致连续的.
\item 证明以上延拓是唯一的, 这是由于任何两个这样的延拓在 $\overline{E}$ 上是连续的映射且在 $\overline{E}$ 的稠密子集 $E$ 上取相同的值.
\end{enumerate}
\end{exercise}

\begin{exercise}\label{exercise:ch09-46}
考虑度量空间的可数族 $\{(X_{n},\rho_{n})\}_{n=1}^{\infty}$. 对这些空间的笛卡儿积 $Z=\prod_{n=1}^{\infty}X_{n}$, 定义 $Z\times Z$ 上的 $\sigma$ 为, 对 $x=\{x_{n}\}$, $y=\{y_{n}\}$,
\begin{align*}
\sigma(x,y)=\sum_{n=1}^{\infty}2^{-n}\rho_{n}^{*}(x_{n},y_{n}),\quad\text{其中每个 }\rho_{n}^{*}=\rho_{n}/(1+\rho_{n}).
\end{align*}
\begin{enumerate}[(i)]
\item 证明 $\sigma$ 是一个度量.
\item 证明 $(Z,\sigma)$ 是完备的当且仅当每个 $(X_{n},\rho_{n})$ 是完备的.
\end{enumerate}
\end{exercise}

\begin{exercise}\label{exercise:ch09-47}
对每个指标 $n$, 定义 $f_{n}(x)=\alpha x^{n}+\beta\cos(x/n)$, $0\leqslant x\leqslant 1$. 参数 $\alpha$ 和 $\beta$ 为何值时序列 $\{f_{n}\}$ 是度量空间 $C[0,1]$ 中的 Cauchy 序列?
\end{exercise}

\begin{exercise}\label{exercise:ch09-48}
令 $\mathcal{D}$ 为由 $(0,1)$ 上可微的连续函数 $f:[0,1]\to\mathbf{R}$ 组成的 $C[0,1]$ 的子空间. $\mathcal{D}$ 是完备的吗?
\end{exercise}

\begin{exercise}\label{exercise:ch09-49}
定义 $\mathcal{L}$ 为由 Lipschitz 函数 $f:[0,1]\to\mathbf{R}$ 组成的 $C[0,1]$ 的子空间. $\mathcal{L}$ 是完备的吗?
\end{exercise}

\begin{exercise}\label{exercise:ch09-50}
对于度量空间 $(X,\rho)$, 根据以下概要完成定理 13 的证明:
\begin{enumerate}[(i)]
\item 若 $\{x_{n}\}$ 和 $\{y_{n}\}$ 是 $X$ 中的 Cauchy 序列, 证明 $\{\rho(x_{n},y_{n})\}$ 是实数的 Cauchy 序列, 且因此收敛.
\item 定义 $X'$ 为 $X$ 中的 Cauchy 序列集. 对于 $X$ 中的两个 Cauchy 序列 $\{x_{n}\}$ 和 $\{y_{n}\}$, 定义 $\rho'(\{x_{n}\},\{y_{n}\})=\lim\rho(x_{n},y_{n})$. 证明这定义了 $X'$ 上的伪度量 $\rho'$.
\item 定义 $X'$ 的两个成员即 $X$ 中的两个 Cauchy 序列 $\{x_{n}\}$ 和 $\{y_{n}\}$ 等价, 若 $\rho'(\{x_{n}\},\{y_{n}\})=0$. 证明这是 $X'$ 中的等价关系且记等价类的集合为 $\widehat{X}$. 定义两个等价类的距离 $\widehat{\rho}$ 为这些类的代表之间的距离 $\rho'$. 证明 $\widehat{\rho}$ 是恰当定义的且是 $\widehat{X}$ 上的一个度量.
\item 证明度量空间 $(\widehat{X},\widehat{\rho})$ 是完备的.
\item 定义从 $X$ 到 $\widehat{X}$ 的映射 $h$ 为对 $x\in X$, $h(x)$ 为所有项是 $x$ 的常数序列的等价类. 证明 $h(X)$ 在 $\widehat{X}$ 中稠密且对所有 $u$, $v\in X$, $\widehat{\rho}(h(u),h(v))=\rho(u,v)$.
\item 定义集合 $\widetilde{X}$ 为 $X$ 与 $\widehat{X}\sim h(X)$ 的不交并. 对 $u$, $v\in\widetilde{X}$, 定义 $\widetilde{\rho}(u,v)$ 如下: 若 $u$, $v\in X$, $\widetilde{\rho}(u,v)=\rho(u,v)$; 若 $u$, $v\in\widehat{X}\sim h(X)$, $\widetilde{\rho}(u,v)=\widehat{\rho}(u,v)$; 而对 $u\in X$, $v\in\widehat{X}\sim h(X)$, $\widetilde{\rho}(u,v)=\widehat{\rho}(h(u),v)$. 从前面两部分得出度量空间 $(\widetilde{X},\widetilde{\rho})$ 是包含 $(X,\rho)$ 作为稠密度子空间的完备度量空间.
\end{enumerate}
\end{exercise}

\begin{exercise}\label{exercise:ch09-51}
证明度量空间 $X$ 的任何两个完备化是等距的. 该等距关系由 $X$ 上的恒等映射给出.
\end{exercise}

\end{document}